\documentclass{article}

\usepackage{arxiv}

\usepackage[utf8]{inputenc} % allow utf-8 input
\usepackage[T1]{fontenc}    % use 8-bit T1 fonts
\usepackage{hyperref}       % hyperlinks
\usepackage{url}            % simple URL typesetting
\usepackage{booktabs}       % professional-quality tables
\usepackage{amsfonts}       % blackboard math symbols
\usepackage{nicefrac}       % compact symbols for 1/2, etc.
\usepackage{microtype}      % microtypography
\usepackage{graphicx}
\usepackage{amsmath}

\title{DFU Segmentation and False Positive Reduction: A Clinical Decision Support System using EfficientNet and Vision Transformers}

\author{
  Alexandre \\
  \texttt{github.com/dfu-segmentation-net} \\
}

\begin{document}
\maketitle

\begin{abstract}
Diabetic Foot Ulcers (DFU) are a severe complication of diabetes, requiring timely and accurate monitoring. While automated segmentation models can assist in wound delineation, they often suffer from high false positive rates when deployed in real-world uncontrolled environments (e.g., confusing socks, floors, or other skin lesions with wounds). This paper presents a robust two-stage deep learning framework designed as a Clinical Decision Support Tool. The first stage utilizes an EfficientNet-B4 U-Net for high-sensitivity binary semantic segmentation. The second stage employs a Vision Transformer (ViT) based Verification Classifier to filter false positives. By integrating the predicted segmentation mask and segmentation percentage as inputs to the ViT, the system effectively distinguishes true wounds from background noise. Our approach achieves a specificity of 93.6\% with the ViT verifier, significantly outperforming baseline CNN methods, and provides a practical workflow for validating image quality before clinical review.
\end{abstract}

\section{Introduction}
Diabetic Foot Ulcers (DFU) represent a significant burden on healthcare systems and patient quality of life. Accurate measurement and monitoring of wound area are critical for assessing healing progress. Manual segmentation is time-consuming and subject to inter-observer variability. Deep learning offers the potential for automated, objective assessment.

However, deploying segmentation models in clinical settings faces a major challenge: \textbf{False Positives}. Models trained on curated wound datasets often over-segment or misclassify non-wound objects (such as socks, furniture, or healthy skin) as ulcers when applied to images taken in uncontrolled environments. In a clinical workflow, excessive false alarms can lead to "alarm fatigue" or distrust in the system.

We propose a Clinical Decision Support System that acts as a "gatekeeper." Instead of replacing the clinician, the system verifies the presence of a wound with high specificity. If the system detects a non-wound or has low certainty, it prompts the user to retake the photo, ensuring that only high-quality, relevant data reaches the medical professional.

\section{Methodology}

\subsection{Dataset Construction}
To build a robust system, we curated a diverse set of data for both training and "hard negative" mining:
\begin{itemize}
    \item \textbf{Positive Samples (Wounds):} The primary training data consists of the Kaggle Wound Segmentation Dataset and the Leprosy Chronic Wound Dataset.
    \item \textbf{Negative Samples (Non-Wounds):} To train the verification classifier, we sourced images from MiniImageNet (general objects), and specific "hard negatives" including Skin Cancer and Skin Disease datasets (e.g., HAM10000). These datasets help the model distinguish DFU from other dermatological conditions.
\end{itemize}

\subsection{Stage 1: Segmentation Network}
The core segmentation model is an \textbf{EfficientNet-B4 U-Net}.
\begin{itemize}
    \item \textbf{Encoder:} EfficientNet-B4 pretrained on ImageNet, selected for its balance of accuracy and computational efficiency.
    \item \textbf{Decoder:} A custom U-Net style decoder with skip connections that bridge the encoder and decoder layers to preserve spatial resolution, critical for boundary delineation.
    \item \textbf{Training:} The model is trained using \textbf{Dice Loss} to handle class imbalance (small wound area vs. large background) and optimized with Adam (LR=0.001).
    \item \textbf{Performance:} This stage is optimized for \textbf{Sensitivity (Recall)}, achieving a Dice Coefficient of approximately 89.3\%. We accept a higher false positive rate at this stage to ensure no potential wounds are missed.
\end{itemize}

\subsection{Stage 2: Verification Classifier}
To address the false positives generated by the high-sensitivity segmentation model, we implement a secondary binary classifier.
\subsubsection{Input Strategy}
We hypothesize that the segmentation map itself contains valuable context (e.g., shape, location) that can aid classification. Therefore, the classifier inputs are:
\begin{enumerate}
    \item \textbf{4-Channel Image:} The original RGB image concatenated with the predicted binary mask from Stage 1.
    \item \textbf{Segmentation Percentage:} A scalar value representing the proportion of the image predicted as "wound." This feature helps filter noise (extremely small regions) or massive failures (entire image segmented).
\end{enumerate}

\subsubsection{Vision Transformer (ViT) Architecture}
We utilized a \textbf{ViT-B/16} architecture for the verification task.
\begin{itemize}
    \item \textbf{Input Modification:} The first convolutional projection layer (\texttt{conv\_proj}) was modified to accept 4 channels. The weights for the mask channel were initialized to zero to allow gradual learning without disrupting the pretrained RGB features.
    \item \textbf{Feature Fusion:} The 768-dimensional feature vector from the ViT encoder is concatenated with the Segmentation Percentage scalar.
    \item \textbf{Classification Head:} A multi-layer perceptron (MLP) maps the combined features to a binary output (Wound vs. Non-Wound).
\end{itemize}

\section{Clinical Workflow and Logic}
The system is designed to function as follows:
\begin{enumerate}
    \item \textbf{Image Capture:} The clinician or patient captures an image of the foot.
    \item \textbf{Segmentation:} The EfficientNet-B4 U-Net generates a candidate mask.
    \item \textbf{Verification:} The ViT model analyzes the image and mask.
    \item \textbf{Decision Logic:}
    \begin{itemize}
        \item If \textbf{Verified as Wound}: The segmentation result and area measurements are presented to the clinician.
        \item If \textbf{Verified as Non-Wound/Low Certainty}: The system triggers a feedback prompt: \textit{"Low certainty detected. Please check lighting, remove background clutter (e.g., socks), and try again."}
    \end{itemize}
\end{enumerate}
We employ \textbf{optimized thresholds} for the verification stage, specifically maximizing specificity to ensure that the "warning" mechanism is accurate and not annoying to the user.

\section{Results}

We evaluated the system on independent test sets, including specific "Specificity Tests" using skin lesion datasets to mimic hard false positives.

\begin{table}[h]
 \caption{Comparison of Verification Models (CNN vs ViT)}
  \centering
  \begin{tabular}{llllll}
    \toprule
    Model & Threshold & Sensitivity & Specificity & F1 Score & Precision \\
    \midrule
    CNN Baseline & 0.5 & 56.81\% & 86.20\% & 60.00\% & 63.68\% \\
    CNN Optimized & 0.41 & \textbf{64.32\%} & 81.80\% & 62.13\% & 60.09\% \\
    ViT Baseline & 0.5 & 61.97\% & 88.60\% & 65.67\% & 69.84\% \\
    \textbf{ViT Optimized} & \textbf{0.70} & 58.69\% & \textbf{93.60\%} & \textbf{67.57\%} & \textbf{79.62\%} \\
    \bottomrule
  \end{tabular}
  \label{tab:results}
\end{table}

\subsection{Analysis}
The ViT model significantly outperforms the CNN baseline in filtering false positives. At the optimized threshold of 0.70, the ViT achieves a \textbf{Specificity of 93.6\%}, reducing the False Positive Rate to just 6.4\% on challenging skin lesion datasets. While there is a slight trade-off in sensitivity compared to the optimized CNN, the substantial gain in precision (+19.53\%) makes the ViT far more suitable for a clinical support tool where trust is paramount.

\section{Conclusion}
We presented a DFU segmentation framework that effectively mitigates false positives through a novel two-stage approach. By leveraging an EfficientNet-B4 U-Net for segmentation and a Vision Transformer that consumes both image and mask data for verification, we achieved high specificity suitable for real-world deployment. This system empowers clinicians by automating wound measurement while actively filtering poor-quality or irrelevant images, streamlining the documentation workflow in diabetic foot care.

\end{document}
